\documentclass[11pt]{amsart}

\usepackage[T1]{fontenc}
\usepackage{lmodern}
\usepackage{amsmath,amssymb,amsthm,mathtools,mathrsfs}
\usepackage[margin=1in]{geometry}
\usepackage{microtype}
\usepackage{enumitem}
\usepackage{float}
\usepackage{xcolor}
\usepackage[colorlinks=true,
  linkcolor=blue!45!black,
  citecolor=blue!45!black,
  urlcolor=blue!45!black]{hyperref}

\newtheorem{theorem}{Theorem}[section]
\newtheorem{proposition}[theorem]{Proposition}
\newtheorem{lemma}[theorem]{Lemma}
\newtheorem{corollary}[theorem]{Corollary}
\theoremstyle{definition}
\newtheorem{definition}[theorem]{Definition}
\theoremstyle{remark}
\newtheorem{remark}[theorem]{Remark}
\newtheorem{warning}[theorem]{Warning}

\DeclareMathOperator{\Dom}{Dom}
\DeclareMathOperator{\Ran}{Ran}
\DeclareMathOperator{\dist}{dist}
\DeclareMathOperator{\Repart}{Re}
\DeclareMathOperator{\Impart}{Im}
\DeclareMathOperator{\spanop}{span}
\newcommand{\C}{\mathbb C}
\newcommand{\R}{\mathbb R}
\newcommand{\Bcal}{\mathcal B}
\newcommand{\Kcal}{\mathcal K}
\newcommand{\Hcal}{\mathcal H}
\newcommand{\Dcal}{\mathcal D}
\newcommand{\eps}{\varepsilon}
\newcommand{\ip}[2]{\left\langle #1,#2\right\rangle}
\newcommand{\norm}[1]{\left\lVert #1\right\rVert}
\newcommand{\abs}[1]{\left\lvert #1\right\rvert}

\title[Conclusions on the Pavlov--Faddeev criterion]
{Conclusions from an Attempt at the\\
Lax--Phillips--Pavlov--Faddeev Regularized Estimate}
\author{OpenAI Codex}
\date{11 August 2026}
\hypersetup{
  pdftitle={Conclusions from an Attempt at the Lax--Phillips--Pavlov--Faddeev Regularized Estimate},
  pdfauthor={OpenAI Codex}
}

\begin{document}
\raggedbottom

\begin{abstract}
We report the outcome of an attempt to prove the regularized
Lax--Phillips--Pavlov--Faddeev estimate associated with the modular
automorphic wave equation.  The physical identifications can be made precise:
the final double-primed interaction semigroup is boundedly similar to an
adjoint compressed-translation model for the pure Blaschke inner function
$\Theta_0(s)=\xi(2s)/\xi(-2s)$, together with a one-dimensional eigenmode at
$-1/2$.  We obtain an exact inverse-Laplace tail formula, an $H^\infty$
quotient formula, a confluent Nevanlinna--Pick criterion retaining all zero
multiplicities, and equivalent completely positive and Lyapunov formulations.
For the actual modular semigroup we prove the unconditional qualitative
decay
\[
 \norm{Z(t)(B+iI)^{-1}}\longrightarrow0,
 \qquad
 -\frac14\leq\omega_{\mathrm{reg}}(B)\leq0.
\]
We also construct pure meromorphic Blaschke models whose entire divisor lies
on $\Re s=1/4$ but whose regularized exponent is $0$.  Hence critical spectral
location, innerness, meromorphicity, divisor symmetry, and root completeness
do not by themselves yield the desired exponent.
Complete positivity does not remove the obstruction: at the required scale it
is exactly equivalent to uniform positivity of all confluent Pick matrices.
The remaining problem is a genuinely uniform, zeta-specific root-synthesis or
weighted-observability estimate.  No unconditional proof or disproof of the
target estimate is obtained, and no claim concerning a proof of the Riemann
hypothesis is made.
\end{abstract}

\maketitle

\tableofcontents

\section{Objective, verdict, and status of possible novelty}

After fixing the positive energy space, the final double-primed
Lax--Phillips interaction semigroup $Z(t)$, and its ordinary $C_0$-generator
$B$, the exact objective was to prove unconditionally
\begin{equation}\label{eq:objective}
 \forall\eps>0\ \exists C_\eps<\infty\ \forall t\geq0:
 \qquad
 \norm{Z(t)(B+iI)^{-1}}
 \leq C_\eps e^{(-1/4+\eps)t}.
\end{equation}
The norm is the operator norm from the positive energy space to itself.  It is
equivalently the norm of $Z(t)$ from $\Dom B$, equipped with
$f\mapsto\norm{(B+iI)f}$, into the energy space.  Neither the
unregularized norm nor a fixed-vector estimate is used below.

The estimate \eqref{eq:objective} was not proved or disproved.  The
identification hypotheses in the governing specification were, however,
verified against the native constructions, and several exact reductions and
obstructions were obtained.  The strongest conclusions are ordered by their
mathematical significance rather than by the order in which they were found.

\begin{remark}[Novelty convention]\label{rem:novelty}
Several elementary consequences below were not located in the sources
inspected during this project.  We call these \emph{candidate observations}
and highlight them explicitly, but make no priority or originality claim.
The main candidates are Theorem~\ref{thm:abstract-counterexample},
Theorem~\ref{thm:qualitative-decay},
Proposition~\ref{prop:finite-dimension}, and
Theorem~\ref{thm:L1-divergence}, together with their zeta-model synthesis.
A dedicated literature search would be required before any novelty claim.
\end{remark}

\section{A critical-line model can have regularized exponent zero}
\label{sec:counterexample}

The following abstract example is the sharpest obstruction found.  It shows
that even placing every model eigenvalue on the critical spectral line does
not control the regularized exponent without quantitative information on root
blocks and their synthesis.

Let $\Theta$ be a pure inner function on $\C_+=\{s:\Repart s>0\}$ and put
\[
 K_\Theta=H^2(\C_+)\ominus\Theta H^2(\C_+),
 \qquad
 S_\Theta(t)=
 \left(P_{K_\Theta}M_{e^{-ts}}\big|_{K_\Theta}\right)^*.
\]
Write $B_\Theta$ for the ordinary generator of $S_\Theta(t)$.

\begin{theorem}[Critical-line Blaschke counterexample]
\label{thm:abstract-counterexample}
There exists a pure meromorphic Blaschke product $\Theta$ such that:
\begin{enumerate}[label=\textup{(\roman*)}]
 \item every zero of $\Theta$ lies on $\Repart a=1/4$, and the divisor is
 invariant under conjugation and under $a\mapsto1/2-\bar a$;
 \item $\Theta$ has zero mean type and no finite boundary singularity;
 \item $S_\Theta(t)$ is a contraction semigroup, $-i\in\rho(B_\Theta)$,
 and the generalized root kernels are complete;
 \item nevertheless,
 \[
  \limsup_{t\to\infty}\frac1t
  \log\norm{S_\Theta(t)(B_\Theta+iI)^{-1}}=0.
 \]
\end{enumerate}
Thus critical-line spectral location, pure innerness, meromorphicity, and the
centered divisor symmetries do not imply the exponent $-1/4$.
\end{theorem}

\begin{proof}
Put $\delta=1/4$, $k_n=n$ for $n\geq2$, and $Y_n=e^{\sqrt n}$.  Let
$\Theta$ be the pure
half-plane Blaschke product whose zeros are
\[
 a_n^\pm=\delta\pm iY_n,
\]
each with multiplicity $k_n+1$.  The Blaschke sum converges because
\[
 \sum_n\frac{2(k_n+1)\delta}{1+\delta^2+Y_n^2}<\infty.
\]
There is no exponential or finite-boundary singular factor by construction,
and there is no finite boundary accumulation.  The divisor has both stated
symmetries.

The zero-kernel derivatives are complete in $K_\Theta$: if
$F\in K_\Theta$ is orthogonal to every derivative kernel, then $F$ vanishes
on the full zero divisor, so $F/\Theta\in H^2(\C_+)$.  Hence
$F\in K_\Theta\cap\Theta H^2=\{0\}$.  The standard model-spectrum theorem
\cite{Nikolski2002} says that its finite generator spectral points are the
zero-derived eigenvalues together with finite boundary singularities.
There are no boundary singularities here and no zero corresponds to $-i$;
therefore $-i\in\rho(B_\Theta)$.

Under the unitary Laplace transform, $S_\Theta(t)$ becomes left translation
on a closed subspace $\mathscr K_\Theta\subset L^2(0,\infty)$ and
$B_\Theta=d/dx$.  Since $a_n^+$ has multiplicity $k_n+1$, the normalized
root vector
\[
 f_n(x)=c_nx^{k_n}e^{-\delta x}e^{iY_nx},
 \qquad
 c_n^2=\frac{(2\delta)^{2k_n+1}}{(2k_n)!},
\]
belongs to $\mathscr K_\Theta\cap W^{1,2}(0,\infty)$ and has norm one.
Direct calculation gives
\begin{equation}\label{eq:counter-graph}
 \norm{(B_\Theta+iI)f_n}_2^2
 =(Y_n+1)^2+\frac{\delta^2}{2k_n-1}.
\end{equation}
The probability density $\abs{f_n(x)}^2dx$ is gamma with shape $2k_n+1$
and rate $2\delta$.  At
\[
 t_n=\frac{2k_n+1}{4\delta},
\]
Chebyshev's inequality yields
\begin{equation}\label{eq:counter-tail}
 \int_{t_n}^\infty\abs{f_n(x)}^2\,dx
 \geq1-\frac4{2k_n+1}.
\end{equation}
Equations \eqref{eq:counter-graph}--\eqref{eq:counter-tail} imply
\[
 \norm{S_\Theta(t_n)(B_\Theta+iI)^{-1}}
 \geq
 \frac{\sqrt{1-4/(2k_n+1)}}
 {\sqrt{(Y_n+1)^2+\delta^2/(2k_n-1)}}.
\]
Since $\log Y_n=\sqrt n=o(t_n)$, the logarithmic limsup is at least zero.
Contractivity and boundedness of the fixed resolvent give the reverse
inequality.

The meromorphic model has no finite boundary spectrum.  None of the stated
qualitative properties changes the calculation above.
\end{proof}

\begin{remark}[Scope of the counterexample]
This is not the zeta inner function and is not a counterexample to RH or to
the Pavlov--Faddeev theorem.  Its purpose is logical: no argument using only
innerness, meromorphic continuation, functional-equation-type divisor
symmetry, completeness, and critical spectral location can prove
\eqref{eq:objective}.  The zeta case must use arithmetic information not
present in the abstract model axioms.
\end{remark}

\section{The strongest unconditional result for the modular semigroup}
\label{sec:qualitative}

The exact Uetake model and its compact-resolvent theorem yield a genuine
operator-norm conclusion, although without a negative exponential rate.

\begin{theorem}[Qualitative regularized decay]
\label{thm:qualitative-decay}
For the final double-primed modular interaction semigroup and its ordinary
generator,
\begin{equation}\label{eq:qualitative-decay}
 \norm{Z(t)(B+iI)^{-1}}\longrightarrow0
 \qquad(t\to\infty).
\end{equation}
Moreover $t\mapsto\norm{Z(t)(B+iI)^{-1}}$ is nonincreasing, and
\begin{equation}\label{eq:omega-bracket}
 -\frac14\leq\omega_{\mathrm{reg}}(B)\leq0.
\end{equation}
\end{theorem}

\begin{proof}
Uetake proves that the reduced generator has meromorphic resolvent, that its
generalized root vectors have dense finite linear span, and that every finite
spectral point is a zeta-derived root eigenvalue; see
\cite[Thms.~4.2 and 4.4]{Uetake2007}.  Every such eigenvalue has strictly
negative real part.  On a root vector the semigroup orbit is a polynomial in
$t$ times $e^{\lambda t}$, hence tends to zero.  Contractivity extends strong
stability from the dense root span to the whole reduced model.  The full
model adds only the mode $e^{-t/2}$, so the physical semigroup is strongly
stable as well.

The generator resolvent is compact by the Lax--Phillips compact-resolvent
theorem, reported in \cite[p.~100]{Uetake2007}; see also
\cite[Thm.~6.17]{LaxPhillips1976}.  Uniformly bounded strong convergence is
uniform on relatively compact sets.  Applying it to the image of the unit
ball under $(B+iI)^{-1}$ proves \eqref{eq:qualitative-decay}.

Since the resolvent commutes with $Z(t)$ and $Z(t)$ is contractive,
\[
 \norm{Z(t+s)(B+iI)^{-1}}
 \leq\norm{Z(t)(B+iI)^{-1}},
\]
which proves monotonicity and the upper bound in
\eqref{eq:omega-bracket}.  A known critical-line zero gives an actual
eigenvalue with real part $-1/4$, and the fixed-resolvent eigenvector test
gives the lower bound.
\end{proof}

\begin{remark}
Norm convergence to zero is strictly stronger than contractive boundedness
and is not a consequence of spectral location alone.  It supplies no
negative exponential exponent: even normal strongly stable semigroups with
compact resolvent can have regularized decay only of order $1/t$.
\end{remark}

\section{Exact scalar reductions: tails, quotient norms, and Pick matrices}
\label{sec:exact-reductions}

Let $\Theta$ be a pure Blaschke inner function on $\C_+$, let
\[
 A_t=P_{K_\Theta}M_{e^{-ts}}\big|_{K_\Theta},
\]
and let $B_\Theta$ generate $A_t^*$.  The results in this section are standard
model-space consequences, but their simultaneous application to the
regularized zeta model is central to locating the exact obstruction.

\begin{proposition}[Graph-tail formula]
\label{prop:tail-formula}
Let $\mathcal L:L^2(0,\infty)\to H^2(\C_+)$ be the unitary Laplace transform
and put $\mathscr K_\Theta=\mathcal L^{-1}K_\Theta$.  Assume
$-i\in\rho(B_\Theta)$.  Then
\[
 \mathcal L^{-1}A_t^*\mathcal L f(x)=f(x+t),
 \qquad
 \Dom B_\Theta=\mathscr K_\Theta\cap W^{1,2}(0,\infty),
 \qquad B_\Theta f=f',
\]
and
\begin{equation}\label{eq:tail-formula}
 \norm{A_t^*(B_\Theta+iI)^{-1}}^2
 =
 \sup_{0\neq f\in\mathscr K_\Theta\cap W^{1,2}}
 \frac{\displaystyle\int_t^\infty\abs{f(x)}^2\,dx}
 {\displaystyle\int_0^\infty\abs{f'(x)+if(x)}^2\,dx}.
\end{equation}
\end{proposition}

\begin{proof}
Multiplication by $e^{-ts}$ corresponds to the right shift under the Laplace
transform, so the adjoint is left translation.  The generator and its domain
follow from the standard shift generator.  Substituting
$h=(B_\Theta+iI)f$ gives \eqref{eq:tail-formula}.
\end{proof}

Assume that $\Theta$ continues holomorphically and nonvanishingly through
$-i$, and set
\begin{equation}\label{eq:r-symbol}
 r(s)=\frac{\Theta(s)/\Theta(-i)-1}{s+i},
 \qquad q_t(s)=e^{-ts}r(s).
\end{equation}
The apparent singularity at $-i$ is removable.  Up to a unimodular scalar,
$r$ is the boundary model kernel $-k_{-i}^{\Theta}$.

\begin{proposition}[$H^\infty$ quotient formula]
\label{prop:quotient-formula}
With the preceding notation,
\begin{equation}\label{eq:quotient-formula}
 \norm{A_t^*(B_\Theta+iI)^{-1}}
 =\norm{P_{K_\Theta}M_{q_t}|_{K_\Theta}}
 =\dist_{H^\infty}(q_t,\Theta H^\infty).
\end{equation}
\end{proposition}

\begin{proof}
Modulo $\Theta H^2$, the identity
\[
 (-s-i)r(s)=1-\frac{\Theta(s)}{\Theta(-i)}
\]
identifies the compressed multiplier by $r$ with the appropriate adjoint
resolvent.  Taking adjoints and commuting the analytic compressed multipliers
gives the first equality.  The second is Sarason's commutant-lifting theorem
\cite{Sarason1967}.
\end{proof}

Let $a_j$ range over the zeros of $\Theta$, with multiplicities $m_j$, and
define
\[
 \psi_t(s)=-\frac{e^{-ts}}{s+i}.
\]
At every zero of $\Theta$, the functions $q_t$ and $\psi_t$ have identical
jets through order $m_j-1$.

\begin{theorem}[Confluent Pick criterion]
\label{thm:confluent-pick}
For $C>0$, the following are equivalent:
\begin{enumerate}[label=\textup{(\roman*)}]
 \item $\norm{A_t^*(B_\Theta+iI)^{-1}}\leq C$;
 \item there exists $F\in H^\infty(\C_+)$ with $\norm F_\infty\leq C$ and
 \[
  F^{(p)}(a_j)=\psi_t^{(p)}(a_j),
  \qquad 0\leq p<m_j;
 \]
 \item for every finite collection of zero blocks, the confluent Pick matrix
 \begin{equation}\label{eq:confluent-pick}
 \left[
 \left.
 \frac1{p!q!}\partial_s^p\partial_{\bar z}^{q}
 \frac{C^2-\psi_t(s)\overline{\psi_t(z)}}{s+\bar z}
 \right|_{s=a_j,z=a_k}
 \right]_{(j,p),(k,q)}
 \end{equation}
 is positive semidefinite.
\end{enumerate}
\end{theorem}

\begin{proof}
Proposition~\ref{prop:quotient-formula} and the pure Blaschke factorization
show that $F-q_t\in\Theta H^\infty$ exactly when the displayed zero jets
agree.  The equivalence with \eqref{eq:confluent-pick} is the confluent
half-plane Nevanlinna--Pick theorem.  Passing from every finite problem to a
global interpolant uses a normal-family diagonal argument.  See
\cite{AglerMcCarthy2000} for the complete Pick framework.
\end{proof}

For simple zeros, \eqref{eq:confluent-pick} reduces to
\begin{equation}\label{eq:simple-pick}
 \left[
  \frac{C^2-w_t(a_j)\overline{w_t(a_k)}}
       {a_j+\bar a_k}
 \right]_{j,k}\succeq0,
 \qquad
 w_t(a)=-\frac{e^{-ta}}{a+i}.
\end{equation}
The one-point minor already forces
\begin{equation}\label{eq:one-point}
 C\geq\frac{e^{-t\Repart a}}{\abs{a+i}}.
\end{equation}
Consequently the desired choice
$C=C_\eps e^{(-1/4+\eps)t}$ for every $\eps>0$ forces
$\Repart a\geq1/4$ for every zero.  In the zeta model the reflected zero
symmetry then forces equality.  Thus even the diagonal Pick conditions
already contain RH; the off-diagonal and confluent conditions contain the
additional nonnormal uniformity.

\section{What complete positivity does and does not provide}
\label{sec:CP}

The Pick matrices admit an exact completely positive formulation.  This
answers the natural operator-algebra question, but also exposes its
circularity at the target scale.

Put
\[
 D_t=A_t^*(B_\Theta+iI)^{-1}.
\]
On a zero kernel $k_a$,
\[
 D_tk_a=\overline{w_t(a)}k_a,
 \qquad w_t(a)=-\frac{e^{-ta}}{a+i},
\]
and on derivative kernels it acts by the corresponding lower-triangular jet
matrix.  Thus Theorem~\ref{thm:confluent-pick} is equivalently the assertion
that this data operator has norm at most $C$ on every finite jet span, with a
bound independent of that span.

\begin{proposition}[Exact UCP reformulation]
\label{prop:UCP}
Assume $K_\Theta$ is infinite dimensional, as in the zeta model.
Let $\widetilde{\Kcal}(K_\Theta)=\Kcal(K_\Theta)+\C I$ be the unitization of
the compact operators.  For $C>0$ define
\[
 \Psi_{t,C}(K+\lambda I)=C^{-2}D_t^*KD_t+\lambda I.
\]
Then
\[
 \Psi_{t,C}\text{ is unital completely positive}
 \quad\Longleftrightarrow\quad
 \norm{D_t}\leq C.
\]
In particular, existence of these UCP maps for
$C=C_\eps e^{(-1/4+\eps)t}$ is exactly the desired regularized estimate.
\end{proposition}

\begin{proof}
If $Q=I-C^{-2}D_t^*D_t\geq0$ and
$\chi(K+\lambda I)=\lambda$, then
\[
 \Psi_{t,C}(A)=C^{-2}D_t^*AD_t+\chi(A)Q
\]
is a sum of completely positive maps and is unital.  Conversely, if
$P_n\uparrow I$ is a finite-rank approximate identity, positivity gives
\[
 0\leq\Psi_{t,C}(I-P_n)=I-C^{-2}D_t^*P_nD_t.
\]
Taking strong limits yields $D_t^*D_t\leq C^2I$.
\end{proof}

Stinespring and Arveson extension therefore do not generate the missing
estimate.  They represent or extend a map after its positivity has been
proved.  In the present problem, subunitality of the one-Kraus map
$X\mapsto D_t^*XD_t$ is precisely the desired norm inequality.  See
\cite{Arveson1969}; the more general operator-argument interpolation theorem
of Muhly--Solel has the same logical form \cite{MuhlySolel2004}.

There is nevertheless an unconditional CP factorization, but only at
exponent zero.  Let $r,q_t$ be as in \eqref{eq:r-symbol} and
$R_0=\norm r_\infty$.  On all of $\C_+$,
\begin{align}\label{eq:CP-zero-rate}
 \frac{R_0^2-q_t(s)\overline{q_t(z)}}{s+\bar z}
={}&
 \frac{R_0^2-r(s)\overline{r(z)}}{s+\bar z}
 \\
 &+r(s)\overline{r(z)}
   \int_0^t e^{-u(s+\bar z)}\,du.\notag
\end{align}
The first summand is a de Branges--Rovnyak kernel and the second is a Gram
kernel.  Hence every confluent restriction is positive and
\[
 \norm{D_t}\leq\norm r_\infty.
\]
Trying to extract $e^{-t/4}$ termwise fails: on the critical line the
remaining phase diagonal is not contractive in the nonorthogonal Cauchy-kernel
Gram metric.  Complete positivity retains positivity but loses exactly the
oscillatory cancellation needed for decay.

\section{Fixed finite root dimension is harmless; uniform dimension is not}
\label{sec:finite-dimension}

The following candidate observation isolates the infinite-dimensional nature
of the obstruction.  It also shows why checking finitely many ordinary Pick
matrices, or only close pairs, cannot settle the estimate.

\begin{proposition}[Uniform estimate at fixed total jet dimension]
\label{prop:finite-dimension}
Let $V\subset\Dom B$ be a subspace of dimension $n$ for a contraction
semigroup, assume $B(V)\subset V$, and suppose
\[
 \sigma(B|_V)\subset\{-x+i\eta:\eta\in\R\}
\]
for some $x>0$.  For every $\eps>0$ there is a constant
$C_{n,\eps,x}$, independent of spectral heights, gaps, and the Jordan
configuration inside $V$, such that
\[
 \norm{e^{tB|_V}(B|_V+iI)^{-1}}
 \leq C_{n,\eps,x}e^{(-x+\eps)t},
 \qquad t\geq0.
\]
Consequently every critical Pick/Hermite test of a fixed total jet dimension
passes with an $\eps$-loss.  The constant is not uniform as $n\to\infty$.
\end{proposition}

\begin{proof}
In an orthonormal Schur basis,
\[
 B|_V=-xI+iD+U,
\]
where $D$ is real diagonal and $U$ is strictly upper triangular.
Dissipativity gives $\abs{U_{jk}}\leq2x$ from each two-dimensional principal
minor of $B+B^*\leq0$, and hence $\norm U\leq Cxn$.  In the interaction
picture $U(s)=e^{-isD}Ue^{isD}$ remains strictly upper triangular.  Every
product of $n$ such matrices vanishes, so the Dyson series terminates:
\[
\norm{e^{tB|_V}}
 \leq e^{-xt}\sum_{k=0}^{n-1}\frac{(Cxn t)^k}{k!}.
\]
Put $D_0=\operatorname{diag}(-x+i(\eta_j+1))$.  Since
$B|_V+iI=D_0+U$ and $D_0^{-1}U$ is strictly upper triangular,
\[
 (B|_V+iI)^{-1}
 =\sum_{k=0}^{n-1}(-D_0^{-1}U)^kD_0^{-1},
 \qquad
 \norm{(B|_V+iI)^{-1}}
 \leq\frac1x\sum_{k=0}^{n-1}(Cn)^k.
\]
Multiply the two estimates and absorb the fixed polynomial into
$e^{\eps t}$.
\end{proof}

For $n=2$ this yields, uniformly as two critical ordinates coalesce and also
for a length-two Jordan block,
\[
 \norm{e^{tB|_V}(B|_V+iI)^{-1}}
 \leq C e^{-xt}(1+t).
\]
Thus pairwise small gaps cause only polynomial transient amplification.  The
actual problem is a bound uniform over arbitrarily large and possibly
nonlocal zero families, together with unbounded jet orders.

\section{Weighted observability, Lyapunov operators, and square functions}
\label{sec:Lyapunov}

Complete positivity becomes useful conceptually when organized as a weighted
observability problem.  It still gives an equivalent certificate rather than
an independent proof.

Let $T(t)$ be a $C_0$ contraction semigroup with generator $B$, and let
$R\in\Bcal(\Hcal)$ commute with $T(t)$.  In the application one may take
$R=(B+iI)^{-1}$ or any other fixed generator resolvent, since fixed resolvent
points give equivalent exponential exponents.

\begin{theorem}[Weighted Gramian and regularized decay]
\label{thm:weighted-Gramian}
Fix $\beta>0$.  The following are equivalent:
\begin{enumerate}[label=\textup{(\roman*)}]
 \item there is $M_\beta<\infty$ such that
 \begin{equation}\label{eq:weighted-observability}
  \int_0^\infty e^{2\beta t}\norm{T(t)Rh}^2\,dt
  \leq M_\beta\norm h^2
  \qquad(h\in\Hcal);
 \end{equation}
 \item there is a bounded positive operator $P_\beta$ satisfying, weakly on
 $\Dom B\times\Dom B$,
 \begin{equation}\label{eq:weak-Lyapunov}
  \ip{P_\beta Bx}{y}+\ip{P_\beta x}{By}
  +2\beta\ip{P_\beta x}{y}
  =-\ip{Rx}{Ry}.
 \end{equation}
\end{enumerate}
Either condition implies
\begin{equation}\label{eq:Gramian-pointwise}
 \norm{T(t)R}\leq C_\beta e^{-\beta t}.
\end{equation}
Conversely, a pointwise estimate with exponent $-\beta'$ for some
$\beta'>\beta$ implies \eqref{eq:weighted-observability}.  Hence the full
family \eqref{eq:objective} is equivalent to existence of such bounded
$P_\beta$ for every $0<\beta<1/4$.
\end{theorem}

\begin{proof}
From \eqref{eq:weighted-observability}, define the bounded positive Gramian
weakly by
\[
 P_\beta=\int_0^\infty e^{2\beta t}
 T(t)^*R^*RT(t)\,dt.
\]
Differentiating its shifted weak form at zero gives
\eqref{eq:weak-Lyapunov}.  Conversely, integrate the derivative of
$e^{2\beta t}\ip{P_\beta T(t)x}{T(t)x}$ and use positivity.

Because $R$ commutes with $T(t)$ and $T(t)$ is contractive, the positive
operators $(T(t)R)^*T(t)R$ decrease in the Loewner order.  Integrating over
$[t-1,t]$ gives
\[
 P_\beta\succeq e^{2\beta t}
 \left(\int_0^1e^{-2\beta u}\,du\right)(T(t)R)^*T(t)R,
\]
which proves \eqref{eq:Gramian-pointwise}.  The converse follows by direct
integration with the strict margin $\beta'>\beta$.  For the epsilon family,
choose $\beta<\beta'<1/4$ in one direction and
$\beta=1/4-\eps$ in the other when $0<\eps<1/4$.  For
$\eps\geq1/4$, contractivity and boundedness of $R$ give the required
nondecaying or growing bound directly.
\end{proof}

The equality \eqref{eq:weak-Lyapunov} is only a quadratic-form statement; no
claim that $P_\beta\Dom B\subset\Dom B^*$ is made.  Similarly,
$X\mapsto T(t)^*XT(t)$ is CP and subunital but need not be norm continuous on
$\Bcal(\Hcal)$.

The outgoing translation representation gives an especially transparent
version.  If $\mathcal T_+\Kcal\subset L^2(\R_-)$ and
$g=\mathcal T_+f$, then
\begin{equation}\label{eq:outgoing-tail}
 \norm{Z(t)f}^2=\int_{-\infty}^{-t}\abs{g(r)}^2\,dr
\end{equation}
and Fubini gives
\begin{equation}\label{eq:outgoing-weight}
 \int_0^\infty e^{2\beta t}\norm{Z(t)f}^2\,dt
 =\frac1{2\beta}\int_{-\infty}^0
 \bigl(e^{-2\beta r}-1\bigr)\abs{g(r)}^2\,dr.
\end{equation}
Thus the missing estimate is precisely boundedness, after the fixed
resolvent, of the exponentially weighted outgoing map
\[
 h\longmapsto e^{-\beta r}\mathcal T_+Rh
 \quad\text{into }L^2(\R_-),
 \qquad \beta<\frac14.
\]
By vector-valued Paley--Wiener, this is equivalently a vertical-line
$L^2$ resolvent estimate for the regularized product
$R(z,B)R$.  It is not an absolute $L^1$ operator-norm contour estimate.

\section{Functional-equation symmetry gives a Krein form, not a positive metric}
\label{sec:detailed-balance}

The zero symmetry $a\mapsto1/2-\bar a$ suggests shifting the generator by
$1/4$ and seeking detailed balance.  The idea fails for a structural reason.

\begin{proposition}[Positive detailed balance annihilates off-critical roots]
\label{prop:DB-annihilates}
Suppose $P\in\Bcal(\Kcal)$ is nonnegative and satisfies the weak identity
\begin{equation}\label{eq:detailed-balance}
 B^*P+PB=-\frac12P,
\end{equation}
equivalently
\[
 Z(t)^*PZ(t)=e^{-t/2}P.
\]
Then $P$ annihilates every root space with
$\Repart\lambda\neq-1/4$.  On a critical Jordan space it annihilates the
nilpotent range.  Consequently an injective positive solution would force all
zeta roots to be critical and semisimple; in the scalar model it would imply
RH together with simplicity of every zero.
\end{proposition}

\begin{proof}
If $Bv=\lambda v$, taking the quadratic form of
\eqref{eq:detailed-balance} on $v$ gives
\[
 \left(2\Repart\lambda+\frac12\right)\ip{Pv}{v}=0.
\]
For an off-critical eigenvalue, positivity gives $Pv=0$.  Induction along a
Jordan chain proves the first assertion.  On a critical block, the induced
nilpotent operator on the quotient by $\ker P$ would be skew-adjoint in a
positive metric, hence zero.
\end{proof}

Bare existence of such a $P$ is uninformative: $P=0$ works, and a known
critical adjoint eigenvector produces a nonzero rank-one solution that simply
annihilates all other dynamics.  A bounded coercive solution would imply the
strictly stronger unregularized estimate
\[
 \norm{Z(t)}\leq C e^{-t/4},
\]
which is not needed for the regularized criterion.

On an off-critical reflected eigenpair, every nondegenerate Hermitian solution
of the centered invariance equation has, in root coordinates, the form
\begin{equation}\label{eq:Krein-pair}
 \begin{pmatrix}0&c\\ \bar c&0\end{pmatrix},
\end{equation}
with signature $(1,1)$.  For paired Jordan blocks the corresponding
anti-diagonal form has signature $(m,m)$.  Thus the functional equation
naturally supplies a Krein pairing whose cross term decays at the centered
rate, not a positive Hilbert metric.  Replacing the form by its absolute value
or square makes it positive but destroys the shifted intertwining relation.

There is a canonical positive substitute in the time model:
\[
 Q_\beta(f,g)=\int_0^\infty e^{2\beta x}f(x)\overline{g(x)}\,dx.
\]
On its natural domain,
\begin{equation}\label{eq:positive-flux-form}
 Q_\beta((B+\beta)f,g)+Q_\beta(f,(B+\beta)g)
 =-f(0)\overline{g(0)}.
\end{equation}
The negative boundary term is outgoing flux.  Graph boundedness of $Q_\beta$
for every $\beta<1/4$ is exactly the weighted formulation of
\eqref{eq:objective}; the endpoint $\beta=1/4$ excludes critical
eigenfunctions and is strictly too strong.  Positivity of the form is
automatic, while its required graph boundedness is the unsolved estimate.

\section{Absolute operator-norm contour integration is impossible}
\label{sec:L1-divergence}

The specification correctly warned that taking the operator norm before a
resolvent contour integral loses too much cancellation.  In the modular
model this warning can be upgraded to a theorem.

\begin{theorem}[Divergence of the absolute resolvent contour]
\label{thm:L1-divergence}
Fix $\eps>0$, put $\alpha=-1/4+\eps$, and let
$\lambda_0\in\rho(B)$.  If the vertical line
$\alpha+i\R$ is contained in $\rho(B)$, then
\begin{equation}\label{eq:L1-divergence}
 \int_{\R}
 \norm{R(\alpha+i\tau,B)^2R(\lambda_0,B)}\,d\tau=\infty.
\end{equation}
If the line meets the spectrum, the proposed contour is already obstructed
by a spectral singularity, which is a pole in the verified meromorphic
modular model.  Thus no proof of \eqref{eq:objective} can be obtained by first
taking the $\Bcal(\Kcal)$-norm and then estimating this contour absolutely.
\end{theorem}

\begin{proof}
Let $\lambda_n=-1/4+i\eta_n$ be a simple critical eigenvalue and let $v_n$
be a corresponding unit eigenvector.  The verified model correspondence
gives such eigenvalues from simple critical zeta zeros.  Applying the three
resolvents to $v_n$ gives
\begin{equation}\label{eq:L1-eigen-lower}
 \norm{R(\alpha+i\tau,B)^2R(\lambda_0,B)}
 \geq
 \frac1{\bigl(\eps^2+(\tau-\eta_n)^2\bigr)
              \abs{\lambda_0-\lambda_n}}.
\end{equation}
An unconditional positive proportion of zeta zeros are simple and lie on the
critical line; Heath--Brown's theorem already suffices
\cite{HeathBrown1979}.  Combining this with the Riemann--von Mangoldt formula
and the standard $O(\log T)$ bound for the number of zeros in an interval of
fixed length, choose one fixed, sufficiently large $q>1$.  In each large
geometric block $[T,qT]$ one may select at least $cT$ such ordinates
$\eta_n=\gamma_n/2$ with a fixed mutual separation.  Trimming the endpoints,
and using alternate blocks if necessary, makes all selected intervals
disjoint.
Integrate \eqref{eq:L1-eigen-lower} over disjoint fixed-radius intervals
centered at the selected ordinates.  Since
$\abs{\lambda_0-\lambda_n}\ll T$ in the block, every interval contributes
$\gg T^{-1}$.  Each geometric block therefore contributes a positive
constant, and summing the blocks proves \eqref{eq:L1-divergence}.
\end{proof}

The failure is specific to absolute $L^1$ control of the operator norm.  It
does not rule out vector-valued $L^2$ contour or square-function methods.  In
fact, Section~\ref{sec:Lyapunov} shows that an $L^2$ estimate with the
regularizer retained is exactly the weighted-observability formulation of
the desired bound.

\section{Verification of the identification chain}
\label{sec:identifications}

This section records which historical and analytic identifications were
actually checked.  They are not consequences of formal similarity of
scattering coefficients; the precise spaces, factors, signs, and
multiplicities matter.

\subsection{The reduced zeta factor is unconditionally pure inner}

We use the entire completed zeta function
\[
 \xi(z)=\frac12z(z-1)\pi^{-z/2}\Gamma(z/2)\zeta(z)
\]
and set
\begin{equation}\label{eq:theta-definition}
 \Theta_0(s)=\frac{\xi(2s)}{\xi(1+2s)}
             =\frac{\xi(2s)}{\xi(-2s)},
 \qquad \Repart s>0.
\end{equation}

\begin{theorem}[Unconditional inner-factor theorem]
\label{thm:inner}
The function $\Theta_0$ belongs to $H^\infty(\C_+)$, has norm one, and is a
meromorphic inner function of zero mean type with no singular inner factor.
Its zeros, with multiplicity, are precisely $a=\rho/2$, where $\rho$ ranges
over the nontrivial zeros of $\xi$.  With factors normalized to equal one at
the origin,
\begin{equation}\label{eq:theta-product}
 \Theta_0(s)=
 \prod_{\rho}
 \frac{\bar\rho}{\rho}\frac{\rho-2s}{2s+\bar\rho}.
\end{equation}
The product is taken with multiplicity, with the zeros ordered by increasing
$\abs{\Impart\rho}$ and conjugate-paired; it converges locally uniformly on
$\C_+$.
Moreover, for the standard Eisenstein coefficient $c$,
\begin{equation}\label{eq:c-theta}
 c\left(\frac12+s\right)
 =\frac{s+1/2}{s-1/2}\Theta_0(s).
\end{equation}
No form of RH is used in these assertions.
\end{theorem}

\begin{proof}
The denominator in \eqref{eq:theta-definition} is nonzero in $\C_+$ because
$\Repart(1+2s)>1$.  Thus the zero divisor and multiplicities follow without
cancellation.  Lagarias's modified-Hadamard-product inequality
\cite[Lemma~2.1]{Lagarias2005} states, for $h\geq1/2$ and
$\Repart u>1/2$, that
\[
 \abs{\xi(h+u)}>\abs{\xi(h+1-\bar u)}.
\]
Taking $h=1/2$ and $u=1/2+2s$ gives
\[
 \abs{\xi(1+2s)}>\abs{\xi(1-2\bar s)}
 =\abs{\xi(2s)},
\]
so $\abs{\Theta_0(s)}<1$ in $\C_+$.  On $s=iy$, the functional equation
and reality symmetry give
$\xi(1+2iy)=\overline{\xi(2iy)}$, with neither side zero.  Hence the boundary
modulus is one at every finite boundary point.

Lagarias is used here only for the Hermite--Biehler inequality and hence the
inner bound.  The remaining factorization assertions follow from the
boundary-meromorphy and asymptotic argument that follows.  The zero counting
formula implies the half-plane Blaschke condition.  Since
$\Theta_0$ is meromorphic across every finite point of $i\R$, the finite
boundary singular measure vanishes.  Direct algebra and Stirling's formula
give, as $x\to+\infty$,
\begin{equation}\label{eq:theta-asymptotic}
 \Theta_0(x)=\sqrt\pi\,x^{-1/2}\bigl(1+O(x^{-1})\bigr),
 \qquad
 \lim_{x\to\infty}\frac1x\log\abs{\Theta_0(x)}=0.
\end{equation}
Thus the exponential factor at infinity is absent.  This proves the pure
Blaschke factorization and \eqref{eq:theta-product}.  Finally, substitution
of the definitions gives
\[
 \Theta_0(s)=
 \frac{s-1/2}{s+1/2}\sqrt\pi\,
 \frac{\Gamma(s)\zeta(2s)}{\Gamma(s+1/2)\zeta(1+2s)},
\]
which is \eqref{eq:c-theta}.
\end{proof}

Here zero mean type means the second limit in
\eqref{eq:theta-asymptotic}.  Some conventions include the exponential mass
at infinity among singular inner factors; the theorem excludes it explicitly.

\subsection{Positive energy and the double-primed Lax--Phillips system}

Let $X=\mathrm{PSL}_2(\mathbb Z)\backslash\mathbb H$, use the positive
Laplacian
\[
 \Delta_X=-y^2(\partial_x^2+\partial_y^2),
\]
and set $A_c=(\Delta_X-1/4I)|_{L_c^2(X)}$.  If $\mathcal E$ denotes the Eisenstein
transform, normalized so that
\[
 \mathcal E A_c\mathcal E^{-1}=M_{\kappa^2},
\]
then on a smooth absolutely continuous core
\begin{equation}\label{eq:energy-map}
 (f_0,f_1)\longmapsto
 \bigl(\kappa\mathcal E f_0,\mathcal E f_1\bigr)
\end{equation}
is isometric for the positive wave energy.  Completion gives the homogeneous
form space
\[
 \dot H^1_{A_c}\oplus L_c^2(X).
\]
The Pavlov--Faddeev absolutely continuous first-order data space and Uetake's
continuous positive-energy space are canonically unitarily equivalent through
this common Eisenstein spectral representation; see
\cite[pp.~164--170]{PavlovFaddeev1975} and
\cite[pp.~102--105]{Uetake2007}.  Thus \eqref{eq:energy-map} supplies the
required energy-unitary identification.  Pavlov--Faddeev use an overall
factor $1/2$ in the quadratic energy, removed by a harmless scalar
normalization.

For the modular surface, the constant mode is the sole below-threshold
residual mode.  Cusp-form standing waves are removed on passing to the
continuous first-order spectral space, and there is no $L^2$ threshold
eigenfunction at $1/4$.  The absence of a threshold eigenvalue is an explicit
assumption in the general Pavlov--Faddeev construction and must not be
silently generalized beyond the modular case.

The correct incoming and outgoing spaces for the orthogonal interaction
semigroup are
\[
 \Dcal_\pm''=\Dcal_\pm^0\cap\Hcal_E',
\]
not the projected primed spaces.  Uetake verifies invariance, trivial
intersection, translate density, and
$\Dcal_-''\perp_E\Dcal_+''$; see
\cite[pp.~104--105]{Uetake2007}.  If
\[
 \Kcal=\Hcal_c'\ominus_E(\Dcal_-''\oplus\Dcal_+''),
\]
then, for $t\geq0$,
\begin{equation}\label{eq:compression}
 Z(t)=P^E_{\Kcal}U(t)|_{\Kcal}
     =Q_+U(t)Q_-|_{\Kcal}
\end{equation}
is a strongly continuous contraction semigroup.  Orthogonality and the
one-sided invariance give \eqref{eq:compression} directly; the discarded
outgoing component remains outgoing, which proves the semigroup law.
Pavlov--Faddeev prove the analogous orthogonal Lax--Phillips axioms and
contraction semigroup for their own incoming and outgoing spaces
\cite[pp.~163--175]{PavlovFaddeev1975}; see also Lax--Phillips
\cite[\S5]{LaxPhillips1980}.

\subsection{The full factor, reduced model, and exceptional mode}

Put
\[
 b_{1/2}(s)=\frac{s-1/2}{s+1/2},
 \qquad \Phi(s)=-b_{1/2}(s)\Theta_0(s).
\]
For the orthogonal double-primed system, the full scattering factor is
\begin{equation}\label{eq:full-scattering}
 S_c''(s)=\Phi(s)=-\frac{s-1/2}{s+1/2}
                     \frac{\xi(2s)}{\xi(-2s)},
\end{equation}
not $c(1/2+s)$ and not $\Theta_0$ alone
\cite[p.~111]{Uetake2007}.  The reciprocal finite factor appearing for
Uetake's nonorthogonal primed spaces is a different convention
\cite[pp.~117--118]{Uetake2007}.
In this subsection write $\Kcal''=\Kcal$ for the full physical
double-primed interaction space.

\begin{theorem}[Exact model and finite exceptional splitting]
\label{thm:exact-model}
Let $\mathcal T_+$ be Uetake's outgoing energy-unitary translation map and
normalize
\[
 [\mathcal L_-g](s)=\int_{-\infty}^0e^{s\tau}g(\tau)\,d\tau,
 \qquad
 \frac1{2\pi}\int_\R\abs{\mathcal L_-g(iy)}^2\,dy=\norm g_2^2.
\]
Then
\begin{equation}\label{eq:U-full}
 [\mathscr U_{\rm full}f](s)
 =[\mathcal L_-\mathcal T_+f](s)
 =\int_0^\infty e^{-sx}(\mathcal T_+f)(-x)\,dx,
 \qquad
 \mathscr U_{\rm full}:\Kcal''\longrightarrow K_\Phi.
\end{equation}
is unitary and intertwines $Z(t)$ with
\[
 A_{\Phi,t}^*
 =\left(P_{K_\Phi}M_{e^{-ts}}|_{K_\Phi}\right)^*.
\]
There is a bounded invariant topological direct sum
\begin{equation}\label{eq:model-splitting}
 K_\Phi=K_{\Theta_0}\dotplus\spanop\{k_{1/2}^{\Phi}\},
\end{equation}
under which
\begin{equation}\label{eq:model-intertwining}
 Z(t)\quad\text{is boundedly similar to}\quad
 A_{0,t}^*\oplus e^{-t/2}.
\end{equation}
More explicitly, if
\[
 J:K_{\Theta_0}\oplus\C\longrightarrow K_\Phi,
 \qquad J(f,c)=f+c\,k_{1/2}^{\Phi},
\]
then the exact physical-to-split-model intertwiner is
\begin{equation}\label{eq:exact-W}
 W=J^{-1}\mathscr U_{\rm full}.
\end{equation}
The reduced generator has spectrum, with algebraic multiplicity,
\begin{equation}\label{eq:model-spectrum}
 \sigma(B_0)=\{-\bar\rho/2:\xi(\rho)=0\}
             =\{-\rho/2:\xi(\rho)=0\},
\end{equation}
its generalized root vectors are complete, and the full generator adds only
the one-dimensional eigenvalue $-1/2$.  There are no other finite spectral
points, and $-i\in\rho(B)$.
\end{theorem}

\begin{proof}
The unitary full-model transform and the generator convention are given in
\cite[pp.~109--113]{Uetake2007}.  Because $b_{1/2}\Theta_0$ has one more
Blaschke zero than $\Theta_0$, $K_{\Theta_0}$ has codimension one in
$K_\Phi$ and is invariant under the adjoint model semigroup.  The reproducing
kernel at $1/2$ is an eigenvector with orbit $e^{-t/2}k_{1/2}^{\Phi}$ and is
not in $K_{\Theta_0}$, since $\Theta_0(1/2)\neq0$.  It therefore supplies the
invariant complement in \eqref{eq:model-splitting}.  The splitting is
bounded because it is a closed codimension-one decomposition, but is not in
general orthogonal.  Thus $J$ is boundedly invertible, and
\eqref{eq:exact-W} proves \eqref{eq:model-intertwining}.

Uetake's Theorems~4.2 and 4.4 identify the reduced resolvent poles, algebraic
multiplicities, and complete root system.  The two forms of the set in
\eqref{eq:model-spectrum} agree by the functional-equation and conjugation
symmetries.  The separate full-space statement
\cite[pp.~100--101]{Uetake2007} gives
$\sigma(B'')\setminus\sigma(B_0)=\{-1/2\}$ and says that the full space is one
dimension larger.  All nontrivial zeta zeros lie in
$0<\Repart\rho<1$, so their generator eigenvalues have real parts strictly
between $-1/2$ and $0$; the exceptional eigenvalue is $-1/2$.  These results
exclude all other finite spectrum and prove $-i\in\rho(B)$.
\end{proof}

The first transform in Theorem~\ref{thm:exact-model} is unitary.  The second,
finite-dimensional splitting is only a bounded similarity.  Hence it gives
two-sided norm comparison and exact equality of logarithmic exponents, not
literal equality of finite-time norms.  The exceptional block decays like
$e^{-t/2}$ and cannot change an exponent at or above $-1/4$.

\subsection{The primary Pavlov--Faddeev theorem}

The primary 1972 paper, in its English translation
\cite[\S\S3--4]{PavlovFaddeev1975}, does prove a whole-space,
resolvent-regularized strip theorem.  Theorem~1 relates bounded holomorphic
continuation of the rationally regularized scattering function to exponential
decay of the uniform interaction-space operator norm
\cite[original pp.~177--180]{PavlovFaddeev1975}; the converse holds in each
strictly smaller strip.  Lemma~8
\cite[pp.~185--186]{PavlovFaddeev1975} supplies the zeta estimates in every
strip of width $a<1/4$ under RH, and the corollary, formulas (48)--(49)
\cite[pp.~186--187]{PavlovFaddeev1975}, states the logarithmic limsup
$-1/4$ with the quantifiers
\[
 \forall a<\frac14\ \exists C_a\ \forall t\geq0.
\]
The paper therefore supports exactly the regularized operator criterion, not
an unregularized norm or a dense fixed-vector substitute.

Its native convention is $Z(t)=e^{itB_{\rm PF}}$.  If $Z(t)=e^{tB}$ in the
ordinary generator convention, then $B=iB_{\rm PF}$ and
\begin{equation}\label{eq:PF-regularizer-conversion}
 (B_{\rm PF}+iI)^{-1}=-i(I-B)^{-1}=-iR(1,B).
\end{equation}
Only after verifying $1,-i\in\rho(B)$ and applying the fixed-resolvent
comparison may this be replaced, at the exponent level, by
$(B+iI)^{-1}=-R(-i,B)$.  Thus the logarithmic exponent and the epsilon family
are unchanged, although the displayed resolvents are not literally equal.
The secondary survey \cite[\S3.2, (3.6), p.~1004]{TakhtajanEtAl2017} quotes
the same regularized formula; the native convention is checked in the primary
paper above.

\section{Other partial results and routes that were tested}
\label{sec:other-routes}

\subsection{The diagonal zero test has stretched-exponential decay}

The fixed regularizer suppresses very high zeros strongly.  The classical
zero-free region, combined with the functional equation, gives
\[
 \Repart a\geq\frac{c}{\log(2+\abs{\Impart a})}
 \qquad(a=\rho/2)
\]
outside a bounded set.  Optimizing in height yields the unconditional
one-point estimate
\begin{equation}\label{eq:diagonal-stretched}
 \sup_{\Theta_0(a)=0}
 \frac{e^{-t\Repart a}}{\abs{a+i}}
 \leq C e^{-c'\sqrt t}.
\end{equation}
This is a useful explanation of the regularizer, but it controls only the
$1\times1$ Pick minors.  It gives no upper bound for the full operator norm,
which depends on arbitrary linear combinations and confluent jets.

\subsection{Ordinary Riesz-basis and Carleson synthesis are unavailable}

On the critical line, the normalized kernel inner product for two nodes
$a=1/4+i\gamma/2$ and $b=1/4+i\gamma'/2$ has modulus
\[
 \frac1{\sqrt{1+(\gamma-\gamma')^2}}.
\]
The Riemann--von Mangoldt density together with the unconditional positive
proportion of simple critical zeros forces arbitrarily close distinct nodes,
so the normalized root kernels are not a Riesz sequence.  More explicitly,
the standard atomic measure
\[
 \mu=\sum_j2\Repart(a_j)\delta_{a_j}
\]
fails the half-plane Carleson condition: unit boundary intervals along a
sequence contain $\gg\log T$ simple critical nodes.  Resolvent weighting may
restore an upper Bessel--Carleson bound, but it supplies neither a lower Riesz
bound nor a bounded coefficient-extraction map.  It therefore does not
identify the graph norm with a weighted $\ell^2$ coefficient norm.

\subsection{Compactness and strong stability have no quantitative rate}

The proof of Theorem~\ref{thm:qualitative-decay} uses all the available
qualitative model information.  It cannot be sharpened abstractly.  For
example, on $\ell^2(\mathbb N)$ let
\[
 Be_n=(-1/n+in)e_n.
\]
Then $e^{tB}$ is a normal, strongly stable contraction semigroup, its
eigenvectors are an orthonormal basis, and $R(1,B)$ is compact, but
\[
 \norm{e^{tB}R(1,B)}\asymp t^{-1}.
\]
Thus even normality, complete roots, compact regularization, and strong
stability give exponent zero without additional quantitative geometry.

\subsection{Why functional-equation pairing cannot be positivized}

There are two distinct reflections.  Scattering gives
\[
 \Theta_0(-s)=\Theta_0(s)^{-1},
\]
which is centered on the Hardy boundary $\Repart s=0$.  Zeta-zero symmetry
gives $a\mapsto1/2-\bar a$, centered on $\Repart a=1/4$.  The latter does not
preserve $\C_+$ or its boundary and does not define a bounded Hardy-space
conjugation.  The canonical model conjugation yields $B\leftrightarrow B^*$
but no $1/4$ shift.  The algebraic critical reflection yields the indefinite
form \eqref{eq:Krein-pair}; applying $Q^*Q$ or $\abs Q$ removes exactly the
off-diagonal relation that encoded the functional equation.

\section{Corrections and qualifications to the governing specification}
\label{sec:specification-errors}

No error was found in the graph-norm equivalence, fixed-resolvent comparison,
formula \eqref{eq:c-theta}, adjoint/time sign, eigenvector lower bound, or the
epsilon-family quantifiers.  The following points should nevertheless be
corrected or sharpened.

\begin{enumerate}[label=\textup{(S\arabic*)},leftmargin=*]
 \item The five labeled identification hypotheses can now be replaced by
 explicit propositions and primary citations, as summarized in
 Section~\ref{sec:identifications}.  They must not remain implicit in a
 claimed proof.

 \item The physical double-primed scattering factor is
 $-b_{1/2}\Theta_0$.  The coefficient $c(1/2+s)=b_{1/2}^{-1}\Theta_0$ has a
 pole at $s=1/2$, whereas the orthogonal double-primed factor has a zero there
 and produces the generator eigenvalue $-1/2$.  These are different objects.

 \item The hypothesis that $\Theta_0$ is bounded inner, of zero mean type,
 and without a singular factor is an unconditional theorem
 (Theorem~\ref{thm:inner}), not an RH-dependent assumption.

 \item In the general Pavlov--Faddeev construction absence of a threshold
 eigenvalue is an explicit assumption.  It is valid for the modular surface
 after the threshold analysis, but not automatically for an arbitrary
 one-cusp group.

 \item Pavlov--Faddeev's energy differs by an overall factor $1/2$ from the
 modern convention.  A scalar rescaling makes the identification literally
 unitary; this does not affect operator norms after both source and target are
 rescaled consistently.

 \item A general bounded similarity provides two-sided constant comparison of
 norms.  Unitarity guarantees exact equality, but it is too strong to say that
 exact equality can occur only under unitarity.

 \item The exceptional spectral-bound argument uses finite dimensionality:
 on the one-dimensional exceptional block the growth exponent equals the
 spectral bound.  Such an inference is false for a general
 infinite-dimensional semigroup.

 \item The phrase zero mean type should be defined, for example by
 \eqref{eq:theta-asymptotic}.  Boundary modulus alone never proves an
 $H^\infty$ bound.

 \item The specification's heuristic warning about absolute contour
 integration can be replaced, under the verified model, by
 Theorem~\ref{thm:L1-divergence}.

 \item Any assertion that the unregularized model norm is exactly one for all
 $t>0$ needs a separate model-space theorem or proof.  It is irrelevant to
 \eqref{eq:objective} and should not be used as a substitute for the
 regularized estimate.
\end{enumerate}

\section{The exact remaining obstruction and acceptance audit}
\label{sec:remaining}

For the zeta inner function, put $a_j=\rho_j/2$, retain the multiplicity
$m_j$, and set $\psi_t(s)=-e^{-ts}/(s+i)$.  After all physical-to-model
identifications and the exceptional splitting, the unresolved assertion is
the following single uniform statement:
\begin{center}
\fbox{\begin{minipage}{0.88\textwidth}
For every $\eps>0$ there is one constant $C_\eps$, independent of $t$, the
number and heights of the selected zeros, their gaps, and their
multiplicities, such that every finite confluent matrix
\begin{equation}\label{eq:boxed-obstruction}
 \left[
 \left.\frac1{p!q!}\partial_s^p\partial_{\bar z}^q
 \frac{C_\eps^2e^{(-1/2+2\eps)t}
       -\psi_t(s)\overline{\psi_t(z)}}{s+\bar z}
 \right|_{s=a_j,z=a_k}
 \right]_{(j,p),(k,q)}
 \succeq0
\end{equation}
for all $t\geq0$ and $0\leq p<m_j$, $0\leq q<m_k$.
\end{minipage}}
\end{center}
Equivalently, one must prove the graph-tail inequality
\begin{equation}\label{eq:remaining-tail}
 \int_t^\infty\abs{f(x)}^2\,dx
 \leq C_\eps^2e^{(-1/2+2\eps)t}
       \int_0^\infty\abs{f'(x)+if(x)}^2\,dx
\end{equation}
for every $f\in\mathscr K_{\Theta_0}\cap W^{1,2}$, with one constant for the
whole graph-norm unit ball.  It is also equivalent to graph boundedness of
all subcritical outgoing weights, or to the bounded positive Lyapunov
operators of Theorem~\ref{thm:weighted-Gramian}.

The one-point part of \eqref{eq:boxed-obstruction} is already RH-equivalent
after zero symmetry.  Assuming RH removes that diagonal obstruction, but it
does not prove the full matrices: Theorem~\ref{thm:abstract-counterexample}
shows that large confluent blocks can destroy the exponent even when every
zero is critical.  Proposition~\ref{prop:finite-dimension} shows the converse
local fact: every fixed total jet dimension is harmless after an
$\eps$-loss.  The unresolved content is precisely uniformity as the total
dimension tends to infinity.

Table~\ref{tab:acceptance} records the audit against the ten acceptance
criteria of the governing specification.  A check means that the issue has
been handled in the reductions above; it does not mean that the final uniform
estimate has been proved.

\begin{table}[H]
\centering
\small
\begin{tabular}{p{0.06\textwidth}p{0.73\textwidth}p{0.08\textwidth}}
\hline
Item & Status & Audit\\
\hline
A1 & Exact physical energy space and physical-to-model transform are stated
     in Section~\ref{sec:identifications}. & check\\
A2 & The positive reduced energy norm is used; the raw indefinite form is
     not used. & check\\
A3 & The exact fixed-resolvent graph norm is retained in
     \eqref{eq:tail-formula}; fixed-resolvent changes are proved harmless. & check\\
A4 & The order $\forall\eps\,\exists C_\eps\,\forall t$ appears explicitly
     in \eqref{eq:objective} and \eqref{eq:boxed-obstruction}. & open\\
A5 & The Pick, tail, and Lyapunov statements all range over the whole
     graph-norm unit ball, not a fixed or dense set of vectors. & open\\
A6 & The full factor $-b_{1/2}\Theta_0$ is split into the reduced model plus
     the faster one-dimensional block. & check\\
A7 & Ordinary generator, time direction, adjoint model, zero map, and bounded
     norm transfer are tracked in Theorem~\ref{thm:exact-model}. & check\\
A8 & Normal-family and weak Gramian limits are justified for the stated
     reductions; absolute contour interchange is rejected by
     Theorem~\ref{thm:L1-divergence}. & qualified\\
A9 & All uniform root-based criteria retain multiplicities through confluent
     jets, and the counterexample treats full Jordan blocks. & check\\
A10 & No growth estimate is inferred from spectrum alone; the abstract
      counterexample proves why this would be invalid. & check\\
\hline
\end{tabular}
\caption{Acceptance audit.  A4--A5 remain open because they are the target
uniform estimate itself.}
\label{tab:acceptance}
\end{table}

\section{Conclusion}

The historical identifications and the analytic inner-function hypothesis
can be discharged rigorously.  The exact problem that remains is not a
spectral-bound calculation: it is the uniform confluent Pick inequality
\eqref{eq:boxed-obstruction}, equivalently the whole-space graph-tail estimate
\eqref{eq:remaining-tail} or a bounded subcritical observability Gramian.
Standard complete positivity, Stinespring dilation, Arveson extension,
scattering unitarity, root completeness, compact resolvent, and the functional
equation each illuminate this statement but do not supply its missing uniform
constant.  The functional equation gives an indefinite critical pairing;
positive forms become useful only after imposing the graph boundedness that
is itself the target.

The strongest unconditional operator conclusion currently obtained is
\[
 \norm{Z(t)(B+iI)^{-1}}\to0,
 \qquad -\frac14\leq\omega_{\rm reg}(B)\leq0.
\]
The critical-line Blaschke construction proves that further abstract model
theory alone cannot close this bracket.  Any successful proof must introduce
new zeta-specific quantitative control of arbitrarily large zero families and
their jets, or an equivalent automorphic observability estimate.  Conversely,
an unconditional proof of that uniform estimate would, already by its
one-dimensional minors and zero symmetry, prove RH.  No such proof or
disproof has been found here.

\begin{thebibliography}{99}

\bibitem{AglerMcCarthy2000}
J.~Agler and J.~E.~McCarthy,
\emph{Complete Nevanlinna--Pick kernels},
J. Funct. Anal. \textbf{175} (2000), no.~1, 111--124,
\href{https://doi.org/10.1006/jfan.2000.3599}{doi:10.1006/jfan.2000.3599}.

\bibitem{Arveson1969}
W.~B.~Arveson,
\emph{Subalgebras of $C^*$-algebras},
Acta Math. \textbf{123} (1969), 141--224,
\href{https://doi.org/10.1007/BF02392388}{doi:10.1007/BF02392388}.

\bibitem{EngelNagel2000}
K.-J. Engel and R.~Nagel,
\emph{One-Parameter Semigroups for Linear Evolution Equations},
Graduate Texts in Mathematics, vol.~194, Springer, New York, 2000.

\bibitem{HeathBrown1979}
D.~R. Heath-Brown,
\emph{Simple zeros of the Riemann zeta-function on the critical line},
Bull. London Math. Soc. \textbf{11} (1979), no.~1, 17--18,
\href{https://doi.org/10.1112/blms/11.1.17}{doi:10.1112/blms/11.1.17}.

\bibitem{Lagarias2005}
J.~C. Lagarias,
\emph{Zero spacing distributions for differenced $L$-functions},
Acta Arith. \textbf{120} (2005), no.~2, 159--184,
\href{https://doi.org/10.4064/aa120-2-4}{doi:10.4064/aa120-2-4}.

\bibitem{LaxPhillips1976}
P.~D. Lax and R.~S. Phillips,
\emph{Scattering Theory for Automorphic Functions},
Annals of Mathematics Studies, vol.~87,
Princeton University Press, Princeton, NJ, 1976.

\bibitem{LaxPhillips1980}
P.~D. Lax and R.~S. Phillips,
\emph{Scattering theory for automorphic functions},
Bull. Amer. Math. Soc. (N.S.) \textbf{2} (1980), no.~2, 261--295,
\href{https://doi.org/10.1090/S0273-0979-1980-14735-7}
{doi:10.1090/S0273-0979-1980-14735-7}.

\bibitem{MuhlySolel2004}
P.~S. Muhly and B.~Solel,
\emph{Hardy algebras, $W^*$-correspondences and interpolation theory},
Math. Ann. \textbf{330} (2004), no.~2, 353--415,
\href{https://doi.org/10.1007/s00208-004-0554-x}
{doi:10.1007/s00208-004-0554-x}.

\bibitem{Nikolski2002}
N.~K. Nikolski,
\emph{Operators, Functions, and Systems: An Easy Reading. Vol.~2:
Model Operators and Systems},
Mathematical Surveys and Monographs, vol.~93,
American Mathematical Society, Providence, RI, 2002.

\bibitem{PavlovFaddeev1975}
B.~S. Pavlov and L.~D. Faddeev,
\emph{Scattering theory and automorphic functions},
J. Soviet Math. \textbf{3} (1975), no.~4, 522--548;
English translation of Zap. Nauchn. Sem. LOMI \textbf{27} (1972), 161--193,
\href{https://www.mathnet.ru/eng/znsl2505}{MathNet source page}.

\bibitem{Sarason1967}
D.~Sarason,
\emph{Generalized interpolation in $H^\infty$},
Trans. Amer. Math. Soc. \textbf{127} (1967), no.~2, 179--203,
\href{https://doi.org/10.1090/S0002-9947-1967-0208383-8}
{doi:10.1090/S0002-9947-1967-0208383-8}.

\bibitem{TakhtajanEtAl2017}
L.~A. Takhtajan, A.~Yu. Alekseev, I.~Ya. Aref'eva,
M.~A. Semenov-Tian-Shansky, E.~K. Sklyanin, F.~A. Smirnov,
and S.~L. Shatashvili,
\emph{Scientific heritage of L. D. Faddeev. Survey of papers},
Russian Math. Surveys \textbf{72} (2017), no.~6, 977--1081,
\href{https://doi.org/10.1070/RM9799}{doi:10.1070/RM9799}.

\bibitem{Uetake2007}
Y.~Uetake,
\emph{The Lax--Phillips infinitesimal generator and the scattering matrix for
automorphic functions},
Ann. Polon. Math. \textbf{92} (2007), no.~2, 99--122,
\href{https://doi.org/10.4064/ap92-2-1}{doi:10.4064/ap92-2-1}.

\end{thebibliography}

\end{document}
